% ==============================================
% Annexe: k_1, k_2, k_3, k_4 Parameters Relation
% Description: Math relation between the k_i
% parameters
% ==============================================

\section{Relación matemática entre los parámetros $k_{i}$}\label{demostralabel}

Partiendo de las Ecuaciones Diferenciales Ordinarias de primer orden \ref{modwp1} y \ref{modwp2} y considerando que estas ecuaciones estan relacionadas entre sí, se puede establecer una EDO de segundo orden que permita expresar una varible en función de otra de la siguiente manera:

% Por una parte sabemos que:
\begin{equation*}
    \frac{dW}{dt} = k_{1}\alpha_{s}W - k_{2}P, \hspace{0.6cm} \frac{dP}{dt} = k_{3}\alpha_{s}W - k_{4}P \hspace{0.2cm}, \hspace{0.3cm} \hspace{0.2cm} \frac{dW}{dt} = \dot{W}, \hspace{0.3cm} \frac{dP}{dt} = \dot{P}
\end{equation*}

% Reescribiendo P en terminos de W obteniendo W**
% \begin{equation*}
%     \frac{dP}{dt} = k_{3}\alpha_{s}W - k_{4}P, \hspace{1cm} \frac{dP}{dt} = \dot{P}
% \end{equation*}

\begin{equation*}
    \ddot{W} = \frac{d}{dt}\left(\frac{dW}{dt}\right) \longrightarrow \ddot{W} = k_{1}\alpha_{s}\frac{dW}{dt} - k_{2}\frac{dP}{dt} \hspace{0.3cm} \longrightarrow \hspace{0.3cm} \ddot{W} = k_{1}\alpha_{s}\dot{W} - k_{2}\dot{P}\hspace{0.15cm}
\end{equation*}

Sabemos que el valor de $\dot{P}$, por lo que reemplazando en $\ddot{W}$ se obtiene:
\begin{equation*}
    \dot{P} = k_{3}\alpha_{s}W - k_{4}P, \hspace{0.2cm}\longrightarrow\hspace{0.2cm} \ddot{W} = k_{1}\alpha_{s}\dot{W} - k_{2}(k_{3}\alpha_{s}W - k_{4}P)
\end{equation*}

% Abrimos parentésis y reescribimos = 0
Separando las variables hacia los lados:
\begin{equation}
    \ddot{W}- k_{1}\alpha_{s}\dot{W} + k_{2}k_{3}\alpha_{s}W = k_{2}k_{4}P
\end{equation}

% suponiendo que k_{2}k_{4}P=0 y cambio de variable
Supóngase  $k_{2}k_{4}P\approx 0$, despejamos $\ddot{W}$, hacemos cambio de variable y resolvemos para $m = \ddot{W}$

% suponiendo que k_{2}k_{4}P=0 y cambio de variable
\begin{equation}
    % (a)m^{2} + (b)m + (c) = 0, \longrightarrow m=\frac{-b \pm \sqrt{b^{2}-4(a)(c)}}{2(a)}
    % \longrightarrow
    m=\left(\frac{-b}{2a}\right) \pm \sqrt{\left(\frac{b}{2a}\right)^{2}-(c)} \hspace{0.2cm}
\end{equation}
Donde $\left[\left(\frac{b}{2a}\right)^{2} - (c)\right]$ decidirá el comportamiento amortiguado del sistema. Reemplazamos valores:
% Reescribimos la cuadrática y vemos la raíz
\begin{equation}
    m=\left[\left(\frac{k_{1}\alpha_{s}}{2}\right)\pm\sqrt{\left(\frac{k_{1}\alpha_{s}^{2}}{2}\right)-(k_{2}k_{3}\alpha_{s})} \hspace{0.2cm}\right] donde \hspace{0.2cm} \left(\frac{k_{1}\alpha_{s}}{2}\right)^{2}\geq k_{2}k_{3}\alpha_{s}
\end{equation}

Despejando la desigualdad, y sabiendo de antemano que
 $\alpha_{s}\neq 0$, se obtiene que:
\begin{equation} \label{relacionkis}
    \left(\frac{k_{1}\alpha_{s}}{2}\right)^{2}\geq k_{2}k_{3}\alpha_{s} \longrightarrow \left(\frac{k_{1}^{2}\alpha_{s}^{2}}{2^{2}}\right)\geq k_{2}k_{3}\alpha_{s} \longrightarrow k_{1}^{2}\geq 4\left(\frac{k_{2}k_{3}}{\alpha_{s}}\right)
\end{equation}
\subsection{Observaciones}
Basado en la representación por EDOs, tanto en las ecuaciones \ref{modwp1} y \ref{modwp2}, así como también por la representación gráfica del modelo, es importante mencionar que, $k_{2}$ debería ser igual a $k_{3}\alpha$, pues parte del peso que no se gana se debe a la producción de producto P que es expulsado del cuerpo del animal, ya sea por la producción de leche L en el proceso de ordeño o por producción de estiércol en el proceso de excreción de las necesidades biológicas del animal.

Por otra parte, el factor $k_{2}k_{4}P$ que se supone aproximadamente igual a 0, puede afectar el comportamiento final del sistema puesto que este valor únicamente sería 0 si el bovino no tuviera perdidas en la producción de desechos.
